\chapter{Modulo FDCScheduling}
\label{ch:fdcscheduling}

\section{Responsabilità}

\module{FDCScheduling} orchestrate la \textbf{generazione e ottimizzazione
degli orari}. Il codice attivo risiede in
\repo{FdC Railway Manager/Services/Scheduling/}; il package SPM
\repo{Packages/FDCScheduling/} è uno scaffold per l'estrazione futura.

\section{Componenti principali}

\begin{table}[htbp]
\centering
\caption{Servizi di scheduling.}
\begin{tabularx}{\textwidth}{lX}
  \toprule
  \textbf{Tipo} & \textbf{Ruolo} \\
  \midrule
  \texttt{ScheduleGenerationEngine} & Orchestratore: crea treni grezzi,
    invoca pipeline, gestisce modalità Takt/libera. \\
  \texttt{ScheduleOptimizationPipeline} & Pipeline a 8 step (vedi \S\ref{sec:pipeline}). \\
  \texttt{TaktEngine} & Allineamento Taktfahrplan, diagnostica hub. \\
  \texttt{ScheduleRefresher} & Ricalcolo orari fisici (arrivi/partenze). \\
  \texttt{ScheduleAIResolver} & Step AI cloud (\texttt{RailwayAIService}). \\
  \texttt{KinematicCalculator} & Tempi di percorrenza su topologia. \\
  \texttt{PathResolver} & Risoluzione percorsi tra stazioni. \\
  \texttt{VehicleSuitabilityEngine} & Compatibilità rotabile/linea. \\
  \bottomrule
\end{tabularx}
\end{table}

\section{Pipeline a 8 step}
\label{sec:pipeline}

\begin{figure}[htbp]
\centering
\begin{tikzpicture}[
  step/.style={draw, rectangle, rounded corners, minimum width=4.8cm,
    minimum height=0.75cm, align=center, font=\small},
  arr/.style={-{Stealth[length=2mm]}}
]
  \node[step, fill=yellow!15] (s1) {1. Shift partenze (greedy)};
  \node[step, fill=blue!10, below=0.5cm of s1] (s2) {2. Refresh fisico};
  \node[step, fill=green!12, below=0.5cm of s2] (s35) {3--5. TaktEngine};
  \node[step, fill=purple!10, below=0.5cm of s35] (s6) {6. AI cloud (opz.)};
  \node[step, fill=orange!12, below=0.5cm of s6] (s7) {7. Ottimizzatore genetico};
  \node[step, fill=red!8, below=0.5cm of s7] (s8) {8. Verifica finale};

  \draw[arr] (s1) -- (s2);
  \draw[arr] (s2) -- (s35);
  \draw[arr] (s35) -- (s6);
  \draw[arr] (s6) -- (s7);
  \draw[arr] (s7) -- (s8);
\end{tikzpicture}
\caption{\texttt{ScheduleOptimizationPipeline.execute}.}
\end{figure}

\subsection{Dettaglio step}

\begin{enumerate}
  \item \textbf{Ottimizzazione partenze} --- Per ogni nuovo treno, prova
    shift coarse ($\pm 10, 20, 30, 50$\,min) e fine ($\pm 1, 2, 5$\,min);
    minimizza i conflitti rilevati da \texttt{ConflictManager}.
  \item \textbf{Refresh fisico} --- \texttt{ScheduleRefresher.refreshMultiple}
    ricalcola tempi di transito e fermata.
  \item \textbf{Allineamento Takt} --- Se almeno un nodo ha
    \texttt{taktMinutes}, \texttt{TaktEngine.generaOrarioCadenzato}
    allinea arrivi/partenze all'hub; GA disabilitato in modalità Takt.
  \item \textbf{AI cloud} --- Opzionale (\texttt{useAI}); delega a
    \texttt{ScheduleAIResolver} e \texttt{RailwayAIService}.
  \item \textbf{GA} --- \texttt{GeneticOptimizer} con 250 iterazioni
    (solo se non Takt).
  \item \textbf{Verifica} --- Secondo refresh + report conflitti residui.
\end{enumerate}

\section{Entry point}

\begin{lstlisting}
// ScheduleGenerationEngine (semplificato)
let topology = RailwayTopology(nodes: network.nodes, edges: network.edges)
let pipeline = ScheduleOptimizationPipeline(topology: topology)
return await pipeline.execute(
    newTrains: trains,
    existingTrains: existing,
    useAI: false,
    useGA: useGA,
    preferredTaktNodeId: taktNode
)
\end{lstlisting}

\section{Dipendenze esterne (non ancora nel package)}

\begin{itemize}
  \item \texttt{ConflictManager} --- Rilevamento conflitti capacità/binari.
  \item \texttt{GeneticOptimizer} --- Meta-euristica su orari.
  \item \texttt{RailwayAIService} --- Ottimizzazione remota.
  \item \texttt{NetworkModel} --- Bridge verso UI e persistenza.
\end{itemize}

L'estrazione in \texttt{FDCScheduling} richiederà protocolli in
\repo{Domain/Services/} (pianificato Fase 5+).

\section{Test di regressione}

\begin{table}[htbp]
\centering
\caption{Suite test scheduling (\texttt{FdC Railway ManagerTests}).}
\begin{tabular}{ll}
  \toprule
  \textbf{File} & \textbf{Copertura} \\
  \midrule
  \texttt{ServicesSchedulingTests} & Pipeline, refresher, path resolver \\
  \texttt{Taktfahrplan120MinTests} & Cicli Takt 120\,min, hub timing \\
  \texttt{TrackAssignmentTests}   & Assegnazione binari \\
  \texttt{EdgeMigrationTests}     & Migrazione modello Edge \\
  \bottomrule
\end{tabular}
\end{table}

\section{Storico}

Il capitolo sostituisce la documentazione del legacy
\texttt{RailwayScheduleOptimizer} (rimosso in commit \texttt{00db4c2}).
Per confronto algoritmico vedere anche
\repo{FdC Railway Manager/RailwayAlgorithmDocs.md}.
